\clearpage
\VerdierPhysicalPageAnchor{0166}%
\begingroup\small\setlength{\parskip}{0pt}\setlength{\emergencystretch}{1.5em}
\VerdierSemanticAnchor{ch4-4-2-3}
\medskip\noindent\textbf{4.2.3.}\quad そこで次のように置く：
\VerdierSemanticAnchor{ch4-4-2-3-1}
\[
E\!\left[\begin{smallmatrix}n\\\omega_3\end{smallmatrix}\right]
=E\!\left[\begin{smallmatrix}n\\f,g,h\end{smallmatrix}\right]
=Z\!\left[\begin{smallmatrix}n\\f,g\end{smallmatrix}\right]
\big/
B\!\left[\begin{smallmatrix}n\\g,h\end{smallmatrix}\right].
\tag{4.2.3.1}
\]
こうして $\mathcal{F}\ell_3(\mathcal J)$ から $\mathcal A$ への関手を
定義した。

いま
$\omega_3=(\mathord\bullet\xrightarrow{f_2}\mathord\bullet
\xrightarrow{f_3}\mathord\bullet\xrightarrow{f_4}\mathord\bullet)$
を $\mathcal{F}\ell_3(\mathcal J)$ の対象とする。
対象
$(\mathord\bullet\xrightarrow{f_3}\mathord\bullet\xrightarrow{f_4}\mathord\bullet)$
は $\mathcal{F}\ell_2(\mathcal J)$ の対象であり、完全列
\[
H^n(f_3)\xrightarrow u H^n(f_4f_3)\longrightarrow H^n(f_4)
\xrightarrow\delta H^{n+1}(f_3).
\]
を与える。さらに、$\mathcal{F}\ell_2(\mathcal J)$ の射：
\VerdierSemanticAnchor{ch4-4-2-3-psi-source}
\[
\begin{tikzcd}[column sep=4em,row sep=2.2em]
{}\arrow[r,"f_2"]\arrow[d,"\operatorname{id}"']&{}
 \arrow[r,"f_3"]\arrow[d,"\operatorname{id}"']&{}\arrow[d,"f_4"]\\
{}\arrow[r,"f_2"']&{}\arrow[r,"f_4f_3"']&{}
\end{tikzcd}
\]
は可換図式
\[
\begin{tikzcd}[column sep=5.5em,row sep=2.7em]
H^n(f_4f_3)\arrow[r,"\delta"]&H^{n+1}(f_2)\\
H^n(f_3)\arrow[u,"u"]\arrow[ur,"\delta"']&{}
\end{tikzcd}
\]
を与える。したがって、図式
\VerdierSemanticAnchor{ch4-4-2-3-psi-diagrams}
\[
\begin{tikzcd}[column sep=4.5em,row sep=2.7em]
H^n(f_4f_3)\arrow[r]\arrow[dr,"\delta"']&
H^n(f_4)\arrow[r,"\delta"]&H^{n+1}(f_3)\\
H^n(f_3)\arrow[u,"u"]\arrow[r,"\delta"']&H^{n+1}(f_2)&{}
\end{tikzcd}
\]
は射
\VerdierSemanticAnchor{ch4-4-2-3-psi}
\[
\Psi^n(\omega_3)\colon
Z\!\left[\begin{smallmatrix}n\\f_3,f_4\end{smallmatrix}\right]
\longrightarrow
H^{n+1}(f_2)\big/
B\!\left[\begin{smallmatrix}n+1\\f_2,f_3\end{smallmatrix}\right].
\]
を定める。
\endgroup
